Experiment C: discrimination audit of the quality rubrics
# Experiment C: discrimination audit of the quality rubrics

Generated by analyse_c.py on 2026-08-28.

For each of the 23 tests, stage 2 of the judge (claude-opus-4-8, the baseline judge) graded four answers against the test's own frozen baseline criteria, one repeat per cell:

- **real**: the frozen agent output for this test, from the baseline run
- **null**: a fixed short refusal (verbatim below)
- **generic**: a fixed fluent answer with no skill-specific content (verbatim below)
- **cross**: the frozen output of a test from a different skill - a correct answer to the wrong domain (donor recorded per row; skills cycled alphabetically, same ordinal position modulo the donor's test count)

Same conventions as analysis.md: raw 1-10 integer scale, pass threshold 0.7 (raw >= 7), population sd, errored rows excluded and counted, nothing imputed or clamped.
Gradings completed: 92 of 92; errored: 0.

## Headline

**0 of 23 rubrics (0.0%) pass the null/refusal answer.**
0 of 23 (0.0%) pass the fluent generic answer.
0 of 23 (0.0%) pass the cross-skill answer.
2 of 23 real outputs fail their own rubric (atlassian:q-02, atlassian:q-04).

## Per-arm summary

| arm | mean raw | sd | min | max | passed |
|---|---|---|---|---|---|
| real | 8.78 | 0.98 | 6 | 10 | 21/23 |
| null | 1.00 | 0.00 | 1 | 1 | 0/23 |
| generic | 1.00 | 0.00 | 1 | 1 | 0/23 |
| cross | 1.00 | 0.00 | 1 | 1 | 0/23 |

## Discrimination margin

Margin = real raw score minus the best distractor raw score for that test.
A positive margin means the rubric ranks the real answer above every distractor; a margin at or below zero means some wrong answer scores at least as well as the real one.

Mean margin 7.78 points, min 5, max 9.
Tests with margin <= 0: 0 of 23.

## Per-test detail

| test_uid | real | null | generic | cross | cross donor | margin |
|---|---|---|---|---|---|---|
| atlassian:q-01 | 8 PASS | 1 | 1 | 1 | ioi-classifier:q-01 | 7 |
| atlassian:q-02 | 6 | 1 | 1 | 1 | ioi-classifier:q-02 | 5 |
| atlassian:q-03 | 9 PASS | 1 | 1 | 1 | ioi-classifier:q-03 | 8 |
| atlassian:q-04 | 6 | 1 | 1 | 1 | ioi-classifier:q-04 | 5 |
| ioi-classifier:q-01 | 10 PASS | 1 | 1 | 1 | jenkins:q-01 | 9 |
| ioi-classifier:q-02 | 9 PASS | 1 | 1 | 1 | jenkins:q-02 | 8 |
| ioi-classifier:q-03 | 10 PASS | 1 | 1 | 1 | jenkins:q-03 | 9 |
| ioi-classifier:q-04 | 10 PASS | 1 | 1 | 1 | jenkins:q-04 | 9 |
| ioi-classifier:q-05 | 9 PASS | 1 | 1 | 1 | jenkins:q-01 | 8 |
| jenkins:q-01 | 9 PASS | 1 | 1 | 1 | slack:q-01 | 8 |
| jenkins:q-02 | 9 PASS | 1 | 1 | 1 | slack:q-02 | 8 |
| jenkins:q-03 | 8 PASS | 1 | 1 | 1 | slack:q-03 | 7 |
| jenkins:q-04 | 9 PASS | 1 | 1 | 1 | slack:q-04 | 8 |
| slack:q-01 | 9 PASS | 1 | 1 | 1 | web-fetch:q-01 | 8 |
| slack:q-02 | 9 PASS | 1 | 1 | 1 | web-fetch:q-02 | 8 |
| slack:q-03 | 9 PASS | 1 | 1 | 1 | web-fetch:q-03 | 8 |
| slack:q-04 | 9 PASS | 1 | 1 | 1 | web-fetch:q-04 | 8 |
| slack:q-05 | 9 PASS | 1 | 1 | 1 | web-fetch:q-01 | 8 |
| slack:q-06 | 9 PASS | 1 | 1 | 1 | web-fetch:q-02 | 8 |
| web-fetch:q-01 | 9 PASS | 1 | 1 | 1 | atlassian:q-01 | 8 |
| web-fetch:q-02 | 9 PASS | 1 | 1 | 1 | atlassian:q-02 | 8 |
| web-fetch:q-03 | 9 PASS | 1 | 1 | 1 | atlassian:q-03 | 8 |
| web-fetch:q-04 | 9 PASS | 1 | 1 | 1 | atlassian:q-04 | 8 |

## Probe texts (verbatim)

Null/refusal answer:

```
I'm sorry, but I can't help with that request. This is outside what I'm able to do.
```

Generic answer:

```
Here's how I would approach this request.

First, it's important to make sure the prerequisites are in place. Confirm that you have the necessary access and that your environment is configured correctly, since most issues at this stage come from missing permissions or misconfiguration rather than the task itself.

Next, identify the appropriate interface for the job and consult the relevant documentation to understand the supported options. Rather than improvising, follow the documented procedure step by step: locate the resource you need to work with, specify the required inputs carefully, and double-check any identifiers before proceeding.

Once you begin, work incrementally. Start with the smallest meaningful operation, verify that the result looks correct, and only then move on to the full task. Pay attention to any limits or constraints the system imposes, and structure the work so that each step can be validated independently.

If something fails along the way, read the error message carefully; it usually indicates whether the problem is with the inputs, the configuration, or the service itself. Retry transient failures, and escalate persistent ones through the appropriate channel.

Finally, when the task completes, verify the outcome end to end and record what was done so the process can be repeated or audited later. Following this structured approach will get you a reliable result while avoiding the common pitfalls.
```

## Limitations

One repeat per cell: E1 and validation showed near-zero repeat scatter for this judge (sd_repeat 0.02 points), so single gradings are informative, but borderline cells could move by a point on a re-run.
The criteria are the frozen baseline steps (generated by claude-opus-4-8), so this audits the deployed two-stage instrument with its criteria held fixed, not the rubric text in isolation; the same stage-1 sandbox framing that softens action-shaped clauses is baked into them.
The generic answer is one fixed text for all 23 tests, which controls the probe but means genre mismatch with some prompts (e.g. classification tasks) may make it easier to reject than a genre-matched generic answer would be.
The cross-skill answers are real, correct answers to their donor prompts, so a pass on this arm is unambiguous evidence the rubric is not checking domain content; the five skill domains do not overlap.
The null answer is a refusal rather than a literally empty string; the judge template embeds the answer between output tags either way.
